% ============================================================
% settings.tex
% Loads all packages and defines every reusable command used
% across sections/. Nothing here changes per job application.
% Only edit this file if you want to change visual formatting.
% ============================================================


% ============================================================
% PACKAGES
% \usepackage{} loads external LaTeX libraries.
% ============================================================

% ------------------------------------------------------------
% geometry — controls page margins.
% All 4 sides set to 0.45in to maximise usable space.
% ------------------------------------------------------------
\usepackage[
  left=0.45in,
  right=0.45in,
  top=0.4in,
  bottom=0.4in
]{geometry}

% ------------------------------------------------------------
% array — extends table column formatting options.
% Needed for the tabular* layouts in entry commands.
% ------------------------------------------------------------
\usepackage{array}

% ------------------------------------------------------------
% ifthen — provides \ifthenelse{}{}{} conditional logic.
% Used by all \IfRole{} commands at the bottom of this file.
% ------------------------------------------------------------
\usepackage{ifthen}

% ------------------------------------------------------------
% enumitem — gives fine control over list spacing.
% Used to tighten bullet point gaps in \ResumeItemize.
% ------------------------------------------------------------
\usepackage{enumitem}

% ------------------------------------------------------------
% hyperref — makes URLs and emails clickable in the PDF.
% hidelinks removes the colored boxes around links.
% ------------------------------------------------------------
\usepackage[hidelinks]{hyperref}

% ------------------------------------------------------------
% xcolor — enables colored text.
% Used by hypersetup below to color links blue.
% ------------------------------------------------------------
\usepackage{xcolor}

% ------------------------------------------------------------
% fontawesome — provides icons like \faGithub, \faEnvelope.
% Used in sections/header.tex for contact icons.
% ------------------------------------------------------------
\usepackage{fontawesome}

% ------------------------------------------------------------
% titlesec — controls section title formatting.
% Used to style \section{} in \ResumeSection below.
% ------------------------------------------------------------
\usepackage{titlesec}

% ------------------------------------------------------------
% tabularx — provides a table that stretches to full width.
% Used in entry commands for left/right aligned rows.
% ------------------------------------------------------------
\usepackage{tabularx}


% ============================================================
% PAGE SETTINGS
% ============================================================

% ------------------------------------------------------------
% pagestyle{empty} removes page numbers and headers/footers.
% A resume should have no page number.
% ------------------------------------------------------------
\pagestyle{empty}

% ------------------------------------------------------------
% Remove all paragraph indentation and spacing between paras.
% Keeps the layout tight and predictable.
% ------------------------------------------------------------
\setlength{\parindent}{0pt}
\setlength{\parskip}{0pt}


% ============================================================
% HYPERLINK COLORS
% colorlinks=true  → links are colored text, not boxed.
% urlcolor=blue    → external URLs appear in blue.
% linkcolor=blue   → internal links appear in blue.
% ============================================================
\hypersetup{
  colorlinks=true,
  urlcolor=blue,
  linkcolor=blue
}


% ============================================================
% REUSABLE COMMANDS
% These are the building blocks used in every section file.
% #1, #2, #3, #4 are positional arguments passed when called.
% ============================================================

% ------------------------------------------------------------
% \ResumeSection{title}
% Renders a bold uppercase section heading with a full-width
% horizontal rule underneath it.
% #1 → section title text (e.g. "Experience")
% ------------------------------------------------------------
\newcommand{\ResumeSection}[1]{%
  \vspace{4pt}%                  % space above the heading
  {\large\bfseries\MakeUppercase{#1}}% large bold uppercase title
  \vspace{1pt}%                  % space between title and rule
  \hrule%                        % full-width horizontal line
  \vspace{4pt}%                  % space below the rule
}

% ------------------------------------------------------------
% \ResumeName{name}
% Renders your name centered at the top in huge bold text.
% #1 → your full name (e.g. "VEDANT GANDHI")
% ------------------------------------------------------------
\newcommand{\ResumeName}[1]{%
  \begin{center}%
    {\Huge\bfseries #1}%         % huge bold name
  \end{center}%
  \vspace{-4pt}%                 % pull next line slightly closer
}

% ------------------------------------------------------------
% \ResumeContact{content}
% Renders contact info centered below the name.
% #1 → raw content (icons, links, separators) from header.tex
% ------------------------------------------------------------
\newcommand{\ResumeContact}[1]{%
  \begin{center}%
    #1%                          % renders whatever is passed in
  \end{center}%
  \vspace{4pt}%                  % space below contact line
}

% ------------------------------------------------------------
% \EducationEntry{institution}{duration}{degree}
% Renders one education entry as a two-column row.
% Left  → bold institution name
% Right → duration (e.g. "2022 -- 2026")
% Below → degree on a full-width second row
% #1 → institution name
% #2 → duration
% #3 → degree
% ------------------------------------------------------------
\newcommand{\EducationEntry}[3]{%
  \noindent%
  \begin{tabular*}{\textwidth}{@{\extracolsep{\fill}}lr}%
    \textbf{#1} & #2\\%          % institution (bold) | duration
    \multicolumn{2}{@{}l}{#3}%   % degree spans full width
  \end{tabular*}%
  \vspace{3pt}%
}

% ------------------------------------------------------------
% \ExperienceEntry{company}{duration}{role}{location}
% Renders one experience entry as a two-column two-row header.
% Row 1 → bold company name | duration
% Row 2 → italic role title  | italic location
% #1 → company name
% #2 → duration (e.g. "Jun 2024 -- Aug 2024")
% #3 → role title
% #4 → location
% ------------------------------------------------------------
\newcommand{\ExperienceEntry}[4]{%
  \noindent%
  \begin{tabular*}{\textwidth}{@{\extracolsep{\fill}}lr}%
    \textbf{#1} & #2\\%          % company (bold) | duration
    \textit{#3} & \textit{#4}%   % role (italic)  | location (italic)
  \end{tabular*}%
  \vspace{2pt}%
}

% ------------------------------------------------------------
% \ProjectEntry{name}{duration}{techstack}
% Renders one project entry header.
% Row 1 → bold project name | duration
% Row 2 → italic tech stack (full width)
% #1 → project name
% #2 → duration (e.g. "2025 -- Present")
% #3 → tech stack (e.g. "React, Node.js, MongoDB")
% ------------------------------------------------------------
\newcommand{\ProjectEntry}[3]{%
  \noindent%
  \begin{tabular*}{\textwidth}{@{\extracolsep{\fill}}lr}%
    \textbf{#1} & #2%            % project name (bold) | duration
  \end{tabular*}\\[-2pt]%        % force tech stack onto its own line
  \textit{#3}%                   % tech stack in italics below
  \vspace{2pt}%
}

% ------------------------------------------------------------
% ResumeItemize environment
% A wrapper around itemize with tightened spacing.
% Use \begin{ResumeItemize} ... \end{ResumeItemize}
% instead of \begin{itemize} for consistent bullet styling.
% ------------------------------------------------------------
\newenvironment{ResumeItemize}
  {\begin{itemize}}%             % opens a standard itemize
  {\end{itemize}}%               % closes it

% Tighten all itemize lists globally (applies to ResumeItemize too).
% leftmargin  → indent of bullets from left edge
% itemsep     → vertical space between bullet items
% topsep      → space before the first bullet
% parsep      → space between paragraphs inside an item
% partopsep   → extra space when list starts a new paragraph
\setlist[itemize]{
  leftmargin=14pt,
  itemsep=1pt,
  topsep=1pt,
  parsep=0pt,
  partopsep=0pt
}

% ------------------------------------------------------------
% \SkillCategory{category}{skills}
% Renders one row of skills with a bold category label.
% #1 → category name (e.g. "Languages")
% #2 → comma-separated skills (e.g. "Java, Python, C++")
% ------------------------------------------------------------
\newcommand{\SkillCategory}[2]{%
  \item \textbf{#1}: #2%
  \vspace{0pt}%
}

% ------------------------------------------------------------
% \Achievement{text}
% Renders one achievement as a bullet item.
% Must be used inside a ResumeItemize environment.
% #1 → achievement text
% ------------------------------------------------------------
\newcommand{\Achievement}[1]{%
  \item #1%      % standard bullet item 
}

% ------------------------------------------------------------
% Spacing helpers — use these instead of raw \vspace{} calls
% to keep spacing consistent and easy to adjust globally.
% ------------------------------------------------------------
\newcommand{\SmallSpace}{\vspace{2pt}}
\newcommand{\MediumSpace}{\vspace{5pt}}


% ============================================================
% ROLE CONDITION COMMANDS
% Each command takes one argument: the content to show.
% Content is only rendered if \Role matches the command's role.
% \ifthenelse{condition}{if true}{if false}
% \equal{a}{b} returns true if a and b are the same string.
%
% Usage example in a section file:
%   \IfBackend{ \WizzerBackend }
%   → only renders \WizzerBackend when \Role is "Backend"
% ============================================================

\newcommand{\IfSDE}[1]{%
  \ifthenelse{\equal{\Role}{SDE}}{#1}{}%
}

\newcommand{\IfBackend}[1]{%
  \ifthenelse{\equal{\Role}{Backend}}{#1}{}%
}

\newcommand{\IfFrontend}[1]{%
  \ifthenelse{\equal{\Role}{Frontend}}{#1}{}%
}

\newcommand{\IfFullStack}[1]{%
  \ifthenelse{\equal{\Role}{FullStack}}{#1}{}%
}

\newcommand{\IfML}[1]{%
  \ifthenelse{\equal{\Role}{ML}}{#1}{}%
}

\newcommand{\IfData}[1]{%
  \ifthenelse{\equal{\Role}{Data}}{#1}{}%
}

\newcommand{\IfDevOps}[1]{%
  \ifthenelse{\equal{\Role}{DevOps}}{#1}{}%
}

\newcommand{\IfQuant}[1]{%
  \ifthenelse{\equal{\Role}{Quant}}{#1}{}%
}

\newcommand{\IfFinance}[1]{%
  \ifthenelse{\equal{\Role}{Finance}}{#1}{}%
}

\newcommand{\IfConsulting}[1]{%
  \ifthenelse{\equal{\Role}{Consulting}}{#1}{}%
}

\newcommand{\IfSRE}[1]{%
  \ifthenelse{\equal{\Role}{SRE}}{#1}{}%
}

\newcommand{\IfMLOps}[1]{%
  \ifthenelse{\equal{\Role}{MLOps}}{#1}{}%
}

\newcommand{\IfCloud}[1]{%
  \ifthenelse{\equal{\Role}{Cloud}}{#1}{}%
}

\newcommand{\IfSecurity}[1]{%
  \ifthenelse{\equal{\Role}{Security}}{#1}{}%
}

\newcommand{\IfApplied}[1]{%
  \ifthenelse{\equal{\Role}{Applied}}{#1}{}%
}

\newcommand{\IfQuantResearch}[1]{%
  \ifthenelse{\equal{\Role}{QuantResearch}}{#1}{}%
}

\newcommand{\IfQuantTrader}[1]{%
  \ifthenelse{\equal{\Role}{QuantTrader}}{#1}{}%
}

\newcommand{\IfAnalytics}[1]{%
  \ifthenelse{\equal{\Role}{Analytics}}{#1}{}%
}

\newcommand{\IfNPCI}[1]{%
  \ifthenelse{\equal{\Role}{NPCI}}{#1}{}%
}

\newcommand{\IfPyDS}[1]{%
  \ifthenelse{\equal{\Role}{PyDS}}{#1}{}%
}


% ============================================================
% ROLE DISPATCH
% Replaces hand-written 18-line \IfRole{} chains in
% sections/experience.tex and sections/projects.tex with a
% single lookup by NAME, using \ifcsname (checks whether a
% command exists without calling it).
%
% Naming convention: for a company/project prefix like
% "Wizzer", the bullet set for each role must be named
% \Wizzer<Role> exactly (e.g. \WizzerSDE, \WizzerBackend).
% Get the name wrong and it's simply treated as "no content" —
% it renders nothing, it does NOT throw a compile error the
% way a bare \IfRole{\SomeTypo} used to.
% ============================================================

% ------------------------------------------------------------
% \RoleContent{prefix}
% Renders \<prefix><Role> if it exists, else renders nothing.
% Use for content where "not relevant to this role" is a
% legitimate outcome (project bullets) — pair with
% \IfRoleContent so the whole entry disappears together.
% ------------------------------------------------------------
\newcommand{\RoleContent}[1]{%
  \expandafter\ifcsname #1\Role\endcsname
    \csname #1\Role\endcsname
  \fi
}

% ------------------------------------------------------------
% \RoleContentOrDefault{prefix}{fallback role name}
% Renders \<prefix><Role> if it exists, else falls back to
% \<prefix><fallback>. Use for real jobs (experience entries),
% which should always show *something* regardless of \Role —
% a real employer should never silently vanish from the resume
% just because that role has no bespoke bullet set yet.
% ------------------------------------------------------------
\newcommand{\RoleContentOrDefault}[2]{%
  \expandafter\ifcsname #1\Role\endcsname
    \csname #1\Role\endcsname
  \else
    \csname #1#2\endcsname
  \fi
}

% ------------------------------------------------------------
% \IfRoleContent{prefix}{block}
% Renders #2 only if \<prefix><Role> exists. Wrap an entire
% project entry (header + tech stack + bullets) in this so the
% header disappears too when there's no content for the
% current role, instead of showing a heading with nothing
% underneath it.
% ------------------------------------------------------------
\newcommand{\IfRoleContent}[2]{%
  \expandafter\ifcsname #1\Role\endcsname
    #2%
  \fi
}

% ------------------------------------------------------------
% \RoleTitle{prefix}{default}
% Looks up \<prefix>Title<Role> for a role-specific job-title
% override (e.g. \AxisMFTitleData = "Data Engineer"). Falls
% back to #2 if no override exists for the current role.
% ------------------------------------------------------------
\newcommand{\RoleTitle}[2]{%
  \expandafter\ifcsname #1Title\Role\endcsname
    \csname #1Title\Role\endcsname
  \else
    #2%
  \fi
}

% ------------------------------------------------------------
% \HideForRole{role}{block}
% Renders #2 for every \Role EXCEPT the one named in #1.
% Use this to suppress an entire experience entry (header +
% bullets) for one specific role, when \RoleContentOrDefault's
% "always show something" behaviour isn't what you want.
%   Example: wrap the whole Techligence \ExperienceEntry block
%   in \HideForRole{SDE}{ ... } to hide it only for SDE, while
%   every other role still shows it normally.
% To hide for more than one role, nest calls:
%   \HideForRole{SDE}{\HideForRole{Backend}{ ... }}
% ------------------------------------------------------------
\newcommand{\HideForRole}[2]{%
  \ifthenelse{\equal{\Role}{#1}}{}{#2}%
}

% ------------------------------------------------------------
% \ProjectSlot{prefix}{block}
% Like \IfRoleContent, but also respects \MaxProjects (set in
% config.tex): once that many projects have already been shown
% for the current \Role, later calls render nothing even if
% content exists for that role. Call sites earlier in
% sections/projects.tex therefore act as higher priority.
% Requires the projectsshown counter (declared below).
% ------------------------------------------------------------
\newcounter{projectsshown}
\setcounter{projectsshown}{0}

\newcommand{\ProjectSlot}[2]{%
  \expandafter\ifcsname #1\Role\endcsname
    \ifnum\value{projectsshown}<\MaxProjects
      #2%
      \stepcounter{projectsshown}%
    \fi
  \fi
}
